\par\noindent\begin{minipage}{\linewidth}
\begin{DefinitionBox}{Reelle Funktion}{real_functions_function}
  Eine \emph{reelle Funktion in einer Variablen}
  \EnglishTerm{real-valued function of one real variable}
  ordnet jedem $x$ aus einer \emph{Definitionsmenge} \EnglishTerm{domain} $D\subseteq\R$ 
  genau einen Wert $y = f(x)$ einer \emph{Zielmenge} \EnglishTerm{codomain} $Z\subseteq\R$ zu.
  \[
    \begin{array}{@{}r@{\quad}l@{}}
      &\begin{aligned}[t]
        f:\ D &\longrightarrow Z,\\
        x &\longmapsto y=\dots
      \end{aligned}\\[1cm]
      \text{oder kurz:}&f(x)=\dots
    \end{array}
    \hspace{1.5cm}
    % Native LaTeX picture inside the formula; deliberately no TikZ.
    \vcenter{\hbox{%
      \setlength{\unitlength}{1cm}%
      \begin{picture}(6.3,3)
        \linethickness{0.6pt}
        \qbezier(1,2.9)(0.15,2.9)(0.15,1.5)
        \qbezier(0.15,1.5)(0.15,0.1)(1,0.1)
        \qbezier(1,0.1)(1.85,0.1)(1.85,1.5)
        \qbezier(1.85,1.5)(1.85,2.9)(1,2.9)
        \qbezier(4.7,2.75)(3.3,2.75)(3.3,1.5)
        \qbezier(3.3,1.5)(3.3,0.25)(4.7,0.25)
        \qbezier(4.7,0.25)(6.1,0.25)(6.1,1.5)
        \qbezier(6.1,1.5)(6.1,2.75)(4.7,2.75)
        \put(1,2.55){\makebox(0,0){$D$}}
        \put(4.7,2.4){\makebox(0,0){$Z$}}
        \put(1,2.1){\makebox(0,0){$\bullet$}}
        \put(1,1.4){\makebox(0,0){$\bullet$}}
        \put(1,0.7){\makebox(0,0){$\bullet$}}
        \put(4.1,2.1){\makebox(0,0){$\bullet$}}
        \put(4.1,1.4){\makebox(0,0){$\bullet$}}
        \put(5.35,1.4){\makebox(0,0){$\bullet$}}
        {\color{THMBlue}\thicklines
          \put(1.12,2.1){\vector(1,0){2.8}}
          \put(1.12,1.42){\vector(1,0){2.8}}
          \put(1.12,0.75){\vector(4,1){2.8}}}
      \end{picture}%
    }}
  \]
\end{DefinitionBox}

Zu einer Funktion gehört ihre Definitionsmenge $D$ und Zielmenge $Z$. Wird in der Kurzschreibweise
nur ein Term $f(x)$ angegeben, geht man von der größtmöglichen reellen Definitionsmenge, 
auf der er $f(x)$ sinnvoll ausgewertet werden kann, aus.

\begin{ExampleBox}{Bestimmung der größtmöglichen Definitionsmenge}{real_functions_domains}
  Für die folgenden drei Terme bestimmen wir jeweils die Bedingungen an $x$
  und daraus die größtmögliche reelle Definitionsmenge.
  \begin{center}
    \renewcommand{\arraystretch}{1.9}
    \setlength{\tabcolsep}{15pt}
    \begin{tabular}{@{}lll@{}}
      \toprule
      Term & Bedingung & Definitionsmenge\\
      \midrule
      $\displaystyle\frac1{x-2}$ & $\Cloze{3cm}{x\ne2}$ &
        $\Cloze{4cm}{\R\setminus\{2\}}$\\
      $\sqrt{x-2}$ & $\Cloze{3cm}{x-2\ge0}$ &
        $\Cloze{4cm}{[2,\infty)}$\\
      $\displaystyle\frac{\sqrt{x-2}}{x-3}$ &
        $\Cloze{3cm}{x\ge2,\ x\ne3}$ &
        $\Cloze{4cm}{[2,3)\cup(3,\infty)}$\\
      \bottomrule
    \end{tabular}
  \end{center}
\end{ExampleBox}
\end{minipage}

\par\noindent\begin{minipage}{\linewidth}
Um den Verlauf einer Funktion sichtbar zu machen, zeichnen wir ihren
\emph{Funktionsgraphen} \EnglishTerm{graph of a function}. Er besteht aus
allen Punkten $(x,f(x))$ mit $x\in D$.
\begin{ExampleBox}{Einen Funktionsgraphen zeichnen}{graph_of_square_function}
  Der Graph der folgenden Funktion heißt \emph{Parabel} \EnglishTerm{parabola}.
  Wir zeichnen ihn mithilfe einer Wertetabelle:
  \[
      \begin{aligned}
        f:\R &\longrightarrow\R,\\
        x &\longmapsto x^2
      \end{aligned}
      \qquad \text{oder} \qquad
      f(x)=x^2
  \]
  \par\smallskip
  \begin{minipage}[c]{0.5\linewidth}
    Wir berechnen einige Funktionswerte und tragen die zugehörigen Punkte
    $(x,f(x))$ in das Koordinatensystem ein:
    \[
      \begin{array}{c||*{4}{w{c}{1.05cm}|}w{c}{1.05cm}}
        x&-2&-1&0&1&2\\ \hline
        f(x)&\Cloze{0.9cm}{4}&\Cloze{0.9cm}{1}&\Cloze{0.9cm}{0}&
          \Cloze{0.9cm}{1}&\Cloze{0.9cm}{4}
      \end{array}
    \]
    Wir tragen die Punkte in das Koordinatensystem ein und vervollstädigen den Graphen.
  \end{minipage}\hfill
  \begin{minipage}[c]{0.46\linewidth}
    \centering
    \CachedTikzFigure{content/resources/chapters/real_functions/tikz/functions_square_graph.tex}
  \end{minipage}
  \par\smallskip
  Gezeigt ist nur ein Ausschnitt: Wegen $D=\R$ setzt sich der Graph beidseitig fort.
\end{ExampleBox}
\end{minipage}

\par\medskip\noindent\begin{minipage}{\linewidth}
Die folgenden beiden Graphen zeigen, warum die eindeutige Zuordnung in der
Funktionsdefinition entscheidend ist.
\par\smallskip
\begin{KeyPointBox}{Eindeutige Zuordnung}{real_functions_unique_assignment}
  Eine Funktion ordnet jedem $x\in D$ genau einen $y$-Wert zu. Daher
  schneidet jede senkrechte Gerade ihren Graphen höchstens einmal.
  \par\smallskip
  \begin{minipage}[t]{0.48\linewidth}
    \centering
    \CachedTikzFigure{content/resources/chapters/real_functions/tikz/functions_unique_graph.tex}
    \par\smallskip
    \raisebox{-0.2ex}{\CachedTikzFigure{content/resources/chapters/real_functions/tikz/functions_assignment_check.tex}}\enspace \textcolor{THMBlue}{Funktion}
    \par
    Jedem $x$ gehört genau ein $y$.
  \end{minipage}\hfill
  \begin{minipage}[t]{0.48\linewidth}
    \centering
    \CachedTikzFigure{content/resources/chapters/real_functions/tikz/functions_nonunique_relation.tex}
    \par\smallskip
    \raisebox{-0.2ex}{\CachedTikzFigure{content/resources/chapters/real_functions/tikz/functions_assignment_cross.tex}}\enspace \textcolor{red!75!black}{Keine Funktion}
    \par
    Einem $x$ entsprechen mehrere $y$-Werte.
  \end{minipage}
\end{KeyPointBox}
\end{minipage}

\par\noindent\begin{minipage}{\linewidth}
Die Parabel im vorigen Beispiel erreicht keine negativen $y$-Werte, obwohl
ihre Zielmenge $Z=\R$ ist. Deshalb betrachten wir die Menge der tatsächlich
angenommenen Werte gesondert.
\begin{DefinitionBox}{\textcolor{green!55!black}{Wertemenge}}{real_functions_range}
  Die \emph{\textcolor{green!55!black}{Wertemenge}} \EnglishTerm{range} $W_f$ einer
  Funktion $f:D\to Z$ enthält genau die Werte, die $f$ tatsächlich annimmt:
  \[
    {\color{green!55!black}W_f}=f(D)=\{f(x)\mid x\in D\}\subseteq Z.
  \]
  \par\noindent\centering
    \CachedTikzFigure{content/resources/chapters/real_functions/tikz/functions_range.tex}
  \par
\end{DefinitionBox}
\begin{ExampleBox}[before skip=0.3\baselineskip]{\textcolor{green!55!black}{Wertemenge} der Quadratfunktion}{real_functions_square_range}
  In Beispiel 5.2 gilt $f:\R\to\R$ mit $f(x)=x^2$.
  Für jedes $x\in\R$ ist $x^2\ge0$; negative Werte treten also nicht auf.
  Umgekehrt wird jedes $y\ge0$ angenommen, denn $f(\sqrt y)=y$.
  Somit enthält die \textcolor{green!55!black}{Wertemenge} genau die nichtnegativen
  reellen Zahlen, während die Zielmenge alle reellen Zahlen umfasst:
  \[
    Z=\R,\qquad
    {\color{green!55!black}W_f=\Cloze{2.5cm}{[0,\infty)}}\subsetneq Z.
  \]
  \par\smallskip\noindent\centering
  \CachedTikzFigure{content/resources/chapters/real_functions/tikz/functions_square_range.tex}
  \par
\end{ExampleBox}
\end{minipage}

\par\noindent\begin{minipage}{\linewidth}
Im Folgenden führen wir weitere Eigenschaften von Funktionen ein und
untersuchen an passenden Funktionsfamilien, was sie aussagen.
\begin{DefinitionBox}{Funktionseigenschaft -- Nullstellen}{real_functions_zeros}
  \begin{minipage}[c]{0.55\linewidth}
    \raggedright
    \emph{Nullstellen} \EnglishTerm{zeros} sind die $x$-Werte aller Schnitt-
    und Berührpunkte des Graphen mit der $x$-Achse. Man berechnet sie aus
    \[
      f(x_i)=0\qquad(x_i\in D).
    \]
    Die zugehörigen Punkte $(x_i,0)$ heißen \emph{$x$-Achsenabschnitte}
    \EnglishTerm{x-intercepts}.
  \end{minipage}\hfill
  \begin{minipage}[c]{0.43\linewidth}
    \centering
    \CachedTikzFigure{content/resources/chapters/real_functions/tikz/functions_zeros.tex}
  \end{minipage}
\end{DefinitionBox}
\end{minipage}
\begin{DefinitionBox}{Funktionseigenschaft -- $y$-Achsenabschnitt}{real_functions_y_intercept}
  \begin{minipage}[c]{0.55\linewidth}
    \raggedright
    Ist $0\in D$, heißt der Funktionswert bei $x=0$
    \emph{$y$-Achsenabschnitt} \EnglishTerm{y-intercept}:
    \[
      y_0=f(0).
    \]
    Der zugehörige Punkt ist $(0,y_0)$. Es gibt höchstens einen
    $y$-Achsenabschnitt.
  \end{minipage}\hfill
  \begin{minipage}[c]{0.43\linewidth}
    \centering
    \CachedTikzFigure{content/resources/chapters/real_functions/tikz/functions_y_intercept.tex}
  \end{minipage}
\end{DefinitionBox}
